\section{Certification of the finite shooting calculation}
\label{cert:appendix}
\label{sec:certificate}
\label{cert:section}

This Appendix proves the finite part of the shooting argument.  The calculation
uses integer interval arithmetic, Taylor expansion on complex disks, and a
separate estimate for the nonlinear dependence on the shooting parameters.
Every terminating decimal in this Appendix denotes its exact rational value.

\begin{proposition}[Finite shooting bounds]\label{cert:shoot}
For every parameter \(z_*\in[-R,R]^2\), where \(R=5\cdot10^{-10}\), let
\((k,d)\) be given by \eqref{prof:parameters}, and let \(p\) be the origin
solution of \eqref{prof:ode} supplied by Lemma~\ref{prof:origin}.
This solution extends to \(0\le s\le S_0=56\) and satisfies
\[
 1+p(s)>0.
\]
For the coordinates of its jet relative to \(p_{\rm a}\), defined in
\eqref{prof:coordinates}, one has
\begin{equation}\label{cert:conclusion}
 |q|\le .25R,\qquad
 \big\|(\Re\xi,\Im\xi)-z_*\big\|_\infty\le .13R.
\end{equation}
\end{proposition}

The proof compares the entire family of shots with a sequence of
reference trajectories at the central parameter \(d=d_0\).  The state
is \((p,p',p'',d)\), with \(d'=0\).  We divide \([1,56]\) into finitely
many blocks.  On each block, the reference trajectory is the exact
solution starting from a rounded three-component jet at its left
endpoint.  At the right endpoint, an interval enclosure supplies the
rounded starting jet for the next block.  Thus successive reference
trajectories may have small jumps at block boundaries; their size is
included in the error estimate.

The interval calculation encloses the reference endpoint jets and
the first-variation matrices of the four-component state across each
block.  Products of these matrices give a linear prediction for the
effect of the two shooting parameters through all blocks.  Rounding the reference jets
and variation matrices introduces separate errors.  A perturbation
estimate compares products of the rounded variation matrices with
the exact transfers and bounds the propagation of those errors.
The origin series supplies the initial central jet, its two parameter
variations, and a bound for the nonlinear parameter remainder.

It remains to compare these linear predictions with the actual
solutions for every \(z_*\in[-R,R]^2\).  The only nonlinear term in the
state equation depends on \(p\); its remainder after linearization
is quadratic in the difference from the reference value.  A uniform
bootstrap bounds this remainder using the certified positive lower
bounds for \(1+p\) along the reference trajectories and the transfer estimates.
It proves continuation and positivity for the whole shooting family and bounds its endpoint
error.  Projecting the endpoint jets into the matching coordinates
then gives \eqref{cert:conclusion}.

We use the maximum norm on vectors and the induced maximum row-sum
norm on matrices, including rectangular matrices.  For a
\(4\times4\) matrix, its superscript \((3)\) denotes its leading
\(3\times3\) block.  We first describe the interval arithmetic and
the series recurrences, then justify their analytic enclosures and
the separate nonlinear comparison.

\subsection{Integer intervals and the two series recurrences}
\label{cert:arithmetic}

Put \(Q_{\rm int}=2^{256}\).  An integer pair \((a,b)\) represents the closed
interval \([a/Q_{\rm int},b/Q_{\rm int}]\).  All arithmetic in
\texttt{verification/certificate.py} is integer arithmetic.  Addition and
multiplication by an integer are exact, with endpoint reversal for a negative
multiplier.  For multiplication, let \(m\) and \(M\) be the minimum and maximum
of the four products of integer endpoints.  The output pair is
\[
 \left(\left\lfloor\frac{m}{Q_{\rm int}}\right\rfloor,
       \left\lceil\frac{M}{Q_{\rm int}}\right\rceil\right).
\]
For a same-sign interval \((a,b)\) excluding zero, the reciprocal pair is
\[
 \left(\left\lfloor\frac{Q_{\rm int}^2}{b}\right\rfloor,
       \left\lceil\frac{Q_{\rm int}^2}{a}\right\rceil\right).
\]
These formulas are valid also for negative intervals.  Division by a positive
integer rounds each endpoint outward in the same way.  Python's integer
division is floor division, so the identity
\(\lceil x\rceil=-\lfloor-x\rfloor\) implements the required upper roundings.
In particular, \texttt{rat(a,b)} encloses the exact rational \(a/b\).
The function \texttt{mag} bounds absolute value; the function \texttt{norm}
bounds the matrix norm specified above.

The operation \texttt{widen} enlarges an interval at each end by the upper
endpoint of a nonnegative error allowance.  The operation \texttt{mid} instead
selects the grid point
\[
 Q_{\rm int}^{-1}\left\lfloor\frac{a+b}{2}\right\rfloor.
\]
For an interval of real width \(w\), its distance from every enclosed number
is at most \(w/2+Q_{\rm int}^{-1}\).  The origin and block estimates below
explicitly charge this error; midpoint selection is never used as an
enclosure operation.

Here are the coefficient identities underlying the computation.
If \(z(\theta)=\sum_{i\ge0}z_i\theta^i\) and
\(b(\theta)=1/z(\theta)=\sum_{i\ge0}b_i\theta^i\), then
\begin{equation}\label{cert:reciprocal}
 b_0=z_0^{-1},\qquad
 b_i=-z_0^{-1}\sum_{t=1}^i z_t b_{i-t}\quad(i\ge1).
\end{equation}
Writing a dot for any one parameter variation gives, for \(i\ge0\),
\begin{equation}\label{cert:tangent-reciprocal}
 \dot b_i=-z_0^{-1}
 \left(\sum_{t=0}^i\dot z_t b_{i-t}
             +\sum_{t=1}^i z_t\dot b_{i-t}\right).
\end{equation}
The second sum is empty at \(i=0\).  In the first sum the coefficient \(b_i\)
has already been computed by \eqref{cert:reciprocal}.  This is exactly the
ordering in \texttt{recip}; it includes the variation of the constant
coefficient \(z_0\).

At the origin, write \(r=s^{3/2}\) locally and
\[
 z=1+p=1+d+kr+\sum_{i\ge2}a_i r^i.
\]
Equation~\eqref{prof:ode} gives
\begin{equation}\label{cert:origin-recurrence}
 a_i=\frac{b_{i-2}-\mathbf1_{\{i=2\}}}
 { (27/8)(i-1)(i+9)(i+10)},\qquad i\ge2,
\end{equation}
where now \(b_i\) are the coefficients of the reciprocal series in \(r\).
The loop in \texttt{origin} computes \(a_2,\ldots,a_{60}\), together with
their two variations under the parameter map \eqref{prof:parameters}.
The weights for evaluation of \(p,p',p''\) at \(s=1\) are
\[
 1,\qquad \frac{3i}{2},\qquad
 \frac{3i}{2}\left(\frac{3i}{2}-1\right).
\]
The code subtracts the constant \(1\) from \(z\) only in the zeroth state
entry, as required.

For an ordinary Taylor step, let \(\ell=1/h\) be the block length,
\(s_c\) the step's initial point, \(\theta=h(s-s_c)\), and
\(\mathfrak r=hs_c+\theta\).  The equation for \(z=1+p\), with
\(E=1+d\), is
\begin{equation}\label{cert:scaled-ode}
 4\mathfrak r^3 z_{\theta\theta\theta}
 +120\mathfrak r^2 z_{\theta\theta}
 +751\mathfrak r z_\theta-1215(z-E)
 =\frac{4}{h^3}\mathfrak r^3(z^{-1}-1).
\end{equation}
This identity also gives the four tangent equations upon differentiation.
For completeness, put \(\mathfrak R=8hs_c\) and
\(F_i=b_i-\mathbf1_{\{i=0\}}\), with negative-index coefficients zero.
The scalar coefficient update used in \texttt{evolve} is
\begin{align}
 4(i+1)(i+2)(i+3)\mathfrak R^3z_{i+3}
 ={}&\frac4{h^3}\sum_{a=0}^3
     c_a\mathfrak R^{3-a}F_{i-a}\notag\\
 &-(96i+960)\mathfrak R^2(i+1)(i+2)z_{i+2}\notag\\
 &-64\mathfrak R(i+1)
       \bigl(12i(i-1)+240i+751\bigr)z_{i+1}\notag\\
 &-512\bigl(4i(i-1)(i-2)+120i(i-1)+751i-1215\bigr)z_i\notag\\
 &-1215\cdot512E\mathbf1_{\{i=0\}},
 \qquad(c_0,c_1,c_2,c_3)=(1,24,192,512).
 \label{cert:step-recurrence}
\end{align}
Indeed \(c_a\mathfrak R^{3-a}\) are the coefficients of
\((\mathfrak R+8\theta)^3\).  The initial three coefficients are
\(z_0,z_\theta(0),z_{\theta\theta}(0)/2\).
For a tangent column one differentiates every term of
\eqref{cert:step-recurrence}; the constant \(E\) has variation \(0,0,0,1\)
in the four columns.  Thus the same update computes both state and tangent
coefficients.  The values \(i=0,\ldots,43\) produce coefficients through
degree \(46\).  Horner evaluation at \(\theta=1/8\), after the appropriate
one or two differentiations, evaluates the three state entries through
degrees \(46,45,44\), respectively.  The analytic justification for their
common remainder allowance is given in Subsection~\ref{cert:disks}.

\subsection{The complete error at the origin}\label{cert:origin-error}

The parameter family of Lemma~\ref{prof:origin} is holomorphic on the
complex polydisc
\[
 |k-k_0|\le .1,\qquad |d-d_0|\le .1,
\]
and its coefficients satisfy
\begin{equation}\label{cert:origin-majorant}
 \sum_{i\ge1}|a_i|3^i<.8,\qquad a_1=k.
\end{equation}
One can check the constants directly from
\eqref{cert:origin-recurrence}: the nonfree denominators exceed \(440\),
\(3|k|<.585\), and \(|1+d|>1.30\).  The majorant improves itself because
\[
 .585+\frac9{440}\left(1+\frac1{1.30-.8}\right)
 =\frac{711}{1100}<.8.
\]
The reciprocal remains nonsingular throughout this construction.
In particular the real solutions are positive up to \(s=1\).
Also \(|d|<.501\), and \eqref{cert:origin-majorant} gives
\[
 |p(1)|<.501+\frac{.8}3<.768,\quad
 |p'(1)|\le .8\max_{i\ge1}\frac{3i/2}{3^i}\le .4,\quad
 |p''(1)|\le .8\max_{i\ge1}
       \frac{(3i/2)(3i/2-1)}{3^i}<.534.
\]
In particular the uniform bound \(2\) for every jet entry is available on
the parameter polydisc.

Let
\begin{equation}\label{cert:origin-tail}
 E_{60}=.8\,\frac{(3\cdot61/2)^2}{3^{61}}
       <5.267\cdot10^{-26}.
\end{equation}
The sequence \((3i/2)^2/3^i\) decreases for \(i\ge61\).
Consequently \(E_{60}\) bounds the omitted tail of each of the three jet
entries.  Cauchy's estimate in either parameter bounds the corresponding
first derivative of this tail at the center by \(10E_{60}\).
For the matrix in \eqref{prof:parameters}, the sum of the absolute values
of all four entries, including the factor \(10^{-8}\), is
\(.01169933\).  Thus the row norm of the omitted tail of the two
\(z_*\)-variation columns is at most \(.12E_{60}\).

Each of the two original parameter displacements has absolute value at most
\(.011R\).  The multivariable Cauchy estimate, with
\(t_*=.11R\), bounds the terms of total degree at least two in the parameter
expansion of any jet entry by
\begin{equation}\label{cert:parameter-tail}
 E_{\rm par}
 =2\sum_{j\ge2}(j+1)t_*^j
 =2\bigl((1-t_*)^{-2}-1-2t_*\bigr)
 <1.816\cdot10^{-20}.
\end{equation}
The factor \(j+1\) counts the monomials of total degree \(j\) in two
parameters.

The width assertion in \texttt{origin} bounds every computed state or
variation interval by \(10^{-20}\).
Let \(U_0\in\R^3\) be the three selected state midpoints.
Let \(Z\) have their two selected variation columns in its first three
rows, and in its fourth row the exact second row of the parameter map
\eqref{prof:parameters}.  The code stores intervals containing that exact
fourth row; its matrix estimates therefore apply to this choice of \(Z\).
Combining the state midpoint error, the state and variation tails,
the two variation midpoint errors, and \eqref{cert:parameter-tail}, gives
\begin{align}
 \left\|
 \begin{pmatrix}p(1)\\p'(1)\\p''(1)\\d\end{pmatrix}
 -\begin{pmatrix}U_0\\d_0\end{pmatrix}-Zz_*
 \right\|_\infty
 &\le 5\cdot10^{-21}+E_{60}+E_{\rm par}
        +.12E_{60}R+10^{-20}R+\frac3{Q_{\rm int}}\notag\\
 &<2.316\cdot10^{-20}<2\cdot10^{-19}.
 \label{cert:initial-error}
\end{align}
The fourth component of the error is exactly zero.
All uncertainty in the interval construction of the central rational
parameters is already included in the width assertion.

\subsection{Analytic continuation and Taylor remainder on a step}
\label{cert:disks}

The 288 blocks run from \(s=1\) to \(s=56\), with the following lengths:
\[
 \begin{array}{c|c|c}
 \text{range}&\ell&\text{number of blocks}\\ \hline
 {[1,2]}&1/32&32\\
 {[2,4]}&1/16&32\\
 {[4,8]}&1/8&32\\
 {[8,56]}&1/4&192.
 \end{array}
\]
Each block consists of eight Taylor steps of length \(\ell/8\).
At every step start the interval assertions imply, for the exact reference
state,
\begin{equation}\label{cert:step-start}
 -.76<p(0)<.4,\qquad |p_\theta(0)|<.06,\qquad
 |p_{\theta\theta}(0)|<.01.
\end{equation}
Every entry in the first three rows of the four-column scaled variation
matrix has modulus less than \(6\); its fourth row is \((0,0,0,1)\).
The step schedule gives \(hs_c\ge32\).

We now justify an analytic bound on the full complex disk
\(|\theta|\le1/2\), using only these starting bounds.
On a disk where \(|p(\theta)-p(0)|\le.1\), one has
\[
 |\mathfrak r|\ge31.5,\qquad |1+p|>.14,\qquad
 |p|<.86,\qquad |p-d_0|<1.3.
\]
Let \(H_2\) be the supremum of \(|p_{\theta\theta}|\) on any smaller
concentric disk.  Integrate the equation for \(p_{\theta\theta}\)
along radial line segments and apply the integral Gronwall inequality.
Since
\[
 \frac{30}{31.5}<.954,\qquad
 \frac{B}{31.5^2}<.190,\qquad
 \frac{|C|}{31.5^3}<.011,\qquad h^{-3}\le.015625,
\]
and \(|p_\theta|\le .06+H_2/2\), this gives
\begin{equation}\label{cert:disk-h2}
 H_2\le
 e^{.954/2}\left[
 .01+\frac12\left\{.190(.06+H_2/2)
               +.011\cdot1.3+.015625\frac{.86}{.14}\right\}\right].
\end{equation}
The coefficient of \(H_2\) on the right is less than one.
Solving this linear inequality yields \(H_2<.123602<.13\).
Therefore
\begin{equation}\label{cert:disk-displacement}
 |p(\theta)-p(0)|\le .06/2+.13/8=.04625<.1.
\end{equation}
This strictly improves the bootstrap.  On every disk of radius at most
\(1/2\), the state and its first two derivatives are bounded and both
denominators \(\mathfrak r\) and \(1+p\) stay uniformly away from zero.
The local holomorphic ODE theorem then continues the solution across any
putative boundary of its disk of existence.  Uniqueness makes the
continuations agree.  Thus the solution is holomorphic on a neighborhood
of the closed disk of radius \(1/2\), with the displayed bounds.

In the scaled state \((p,p_\theta,p_{\theta\theta},d)\), the Jacobian is
\begin{equation}\label{cert:scaled-jacobian}
 \mathcal J(\theta)=
 \begin{pmatrix}
 0&1&0&0\\
 0&0&1&0\\
 -C/\mathfrak r^3-h^{-3}/(1+p)^2&
       -B/\mathfrak r^2&-30/\mathfrak r&C/\mathfrak r^3\\
 0&0&0&0
 \end{pmatrix}.
\end{equation}
On the complex disk its norm is bounded by the maximum of \(1\) and
\[
 \frac{30}{31.5}+\frac{751/4}{31.5^2}
       +\frac{1215/2}{31.5^3}+\frac{1/64}{.14^2}<2.1.
\]
Each actual variation column, with initial norm at most \(6\), consequently
has norm at most \(6e^{2.1/2}<17.146<32\) throughout the disk.
The state entries themselves also have modulus less than \(32\).
These estimates apply to the actual block variation already propagated
from the block's initial identity matrix; a new identity is not assumed
at each substep.

If a scalar holomorphic function is bounded by \(32\) on this disk,
Cauchy's coefficient estimate gives a tail after degree \(44\), at
\(\theta=1/8\), bounded by
\begin{equation}\label{cert:step-tail}
 32\sum_{i\ge45}\left(\frac{1/8}{1/2}\right)^i
 =\frac{128}{3}\,4^{-45}<48\,4^{-45}.
\end{equation}
Every state entry and every tangent entry is approximated to at least
degree \(44\).  Hence the allowance \(48\,4^{-45}\) added by
\texttt{evolve} encloses all their Taylor tails.  The interval recurrences
then inductively enclose the exact reference solution and its exact
first variation through all eight substeps.

On a real substep, the sharper displacement bound is
\begin{equation}\label{cert:real-step-displacement}
 |p(\theta)-p(0)|
 \le .06/8+\frac{.13}{2\cdot8^2}
 =.008515625<.01.
\end{equation}
Accordingly, the smallest lower starting endpoint for \(1+p\), reduced
by \((\lfloor Q_{\rm int}/100\rfloor+1)/Q_{\rm int}>.01\), is a valid
lower bound \(m_j\) on the entire reference block.  This is the variable
\texttt{mini} divided by \(Q_{\rm int}\).

\subsection{The finite inequalities and their exact interpretation}
\label{cert:ledger}

Write the block endpoints as \(s_0=1,\ldots,s_N=56\), \(N=288\).
On block \(j\), the reference solution has initial physical jet \(U_j\)
and parameter \(d_0\); it is evolved exactly on that block.
The rounded jet selected for the next block is \(U_{j+1}\).
Let \(M_j\) be the exact \(4\times4\) transfer of the linearized
four-component physical state along this reference block.

The conversion from scaled derivatives to physical derivatives is essential.
With
\[
 S_j=\diag(1,\ell_j,\ell_j^2,1),
\]
a scaled transfer \(T\) becomes \(S_j^{-1}TS_j\).
Thus its entry in row \(a\), column \(b\), for \(a,b\in\{0,1,2\}\),
is multiplied by \(h_j^a/h_j^b\); column \(3\) is multiplied by \(h_j^a\).
The fourth row remains \((0,0,0,1)\).
The code performs this conversion before asserting the widths.
Consequently its assertions imply
\begin{equation}\label{cert:local-errors}
 \big\|U_{j+1}- (p_{\rm ref},p'_{\rm ref},p''_{\rm ref})
                         (s_{j+1})\big\|_\infty<10^{-19},
 \qquad
 |(M_j-\widetilde M_j)_{ab}|<10^{-15},
\end{equation}
where \(\widetilde M_j\) is the selected matrix of transfer midpoints.
The state and transfer widths are, respectively, less than
\(10^{-19}\) and \(10^{-15}\); the stated looser midpoint bounds follow
from the midpoint estimate in Subsection~\ref{cert:arithmetic}.
All fourth-component errors in \eqref{cert:local-errors} are zero.

Define the products of these point matrices by
\[
 \widetilde T_{i,j}=
 \widetilde M_{i-1}\cdots\widetilde M_j\quad (i>j),
 \qquad \widetilde T_{i,i}=I_4.
\]
The exact integer comparisons in \texttt{check} give
\begin{align}
 m_j&>.23,\qquad \|Z\|<.02,\notag\\
 \max_{i\ge j}\|\widetilde T_{i,j}\|&\le8500,
 &\max_i\sum_{j<i}\|\widetilde T_{i,j+1}^{(3)}\|&\le570000,\notag\\
 \max_i\sum_{j<i}\frac{\ell_j}{m_j^3}
                \|\widetilde T_{i,j+1}^{(3)}\|&\le246000,
 &\max_i\|\widetilde T_{i,0}Z\|&\le2.7.
 \label{cert:finite-products}
\end{align}
Here and below a maximum over endpoints includes endpoint \(0\);
the assertion for it follows from \(\|Z\|<.02\).

To explain the summation indices, for each endpoint \(i\) the inner loop
starts with the identity and takes \(j=i-1,\ldots,0\).
Immediately before multiplying by \(\widetilde M_j\), its current interval
matrix encloses \(\widetilde T_{i,j+1}\).
The update \texttt{T=mm(T,allm[j])} therefore produces the desired forward
product in the correct order.  If \(r_j=Q_{\rm int}m_j\) and \(\mathfrak d\)
is the integer interval bound for the leading \(3\times3\) norm, then
\[
 \frac{\mathfrak d Q_{\rm int}^{\,2}}{h_jr_j^3}
 =\frac{\ell_j}{m_j^3}\frac{\mathfrak d}{Q_{\rm int}}.
\]
The expression used for \texttt{K} is its floor plus one, hence an upper
bound even when the quotient is integral.  Thus the interval product
arithmetic and the weighted sum both round in the required direction.

Let \(P_{\rm a}(d)=(p_{\rm a},p'_{\rm a},p''_{\rm a})(56)\), and put
\(Y=(\widetilde T_{N,0}Z)_{\mathrm{top}\,3}\).
The final comparisons in \texttt{check} give
\begin{align}
 \|U_N-P_{\rm a}(d_0)\|_\infty&<.04R,&
 \|\Lambda(U_N-P_{\rm a}(d_0))\|_\infty&<.006R,\notag\\
 \|(\Lambda Y)_{1,\cdot}\|_{\rm row}&<.116,&
 \|(\Lambda Y)_{\mathrm{bottom}\,2}-I_2\|&<.001.
 \label{cert:finite-matching}
\end{align}
The jet factors used for the term \(d_m s^{-3m}\) are
\(1,-3m,3m(3m+1)\), with denominators \(56^{3m}\),
\(56^{3m+1}\), \(56^{3m+2}\), respectively.
After multiplication by \(d_m\), these give its value and first two
derivatives at \(s=56\).
The index test in the last assertion subtracts the \(2\times2\) identity
from the two nonreal-coordinate rows.

The listed comparisons are finite integer statements: the displayed source
constructs their operands and asserts them.  Subsection~\ref{cert:provenance} gives the command for checking them
with assertions enabled.
The preceding recurrence and disk arguments prove that these finite
statements bound the specified analytic trajectories and variations.

\subsection{Transfer errors and estimates within a block}
\label{cert:transfer-errors}

Let \(T_{i,j}=M_{i-1}\cdots M_j\) be the exact products.
By \eqref{cert:local-errors},
\(\|M_j-\widetilde M_j\|<4\cdot10^{-15}\).
Expand a product of at most \(288\) factors, choosing the error factor at
\(k\) positions.  Each intervening, possibly empty, product of approximate
factors has norm at most \(8500\), by \eqref{cert:finite-products}.
Summing the bounds for \(k\ge1\) gives
\begin{equation}\label{cert:product-perturbation}
 \|T_{i,j}-\widetilde T_{i,j}\|
 \le8500\left[(1+4\cdot10^{-15}\cdot8500)^{288}-1\right]
 <.001.
\end{equation}
In particular \(\|T_{i,j}\|<8501\).  Since the lengths of all the
blocks sum to \(55\), \eqref{cert:finite-products} implies
\begin{equation}\label{cert:exact-sums}
 \max_i\sum_{j<i}\|T_{i,j+1}^{(3)}\|\le580000,\qquad
 \max_i\sum_{j<i}\frac{\ell_j}{m_j^3}
             \|T_{i,j+1}^{(3)}\|\le250000,
\end{equation}
using the elementary comparisons
\[
 570000+288\cdot.001<580000,\qquad
 246000+\frac{55\cdot.001}{.23^3}<250000.
\]
Also
\begin{equation}\label{cert:linear-size}
 \|T_{i,0}Z\|\le2.7+.001\cdot.02<2.72.
\end{equation}

We need two estimates for transfers between any two ordered points of
one reference block, including points in different Taylor substeps.
The real scaled Jacobian \eqref{cert:scaled-jacobian} has norm at most
\[
 L_{\rm real}=
 \max\left\{1,\frac{30}{32}+\frac{751/4}{32^2}
                  +\frac{1215/2}{32^3}
                  +\frac{1/64}{.23^2}\right\}.
\]
An interval of at most one block unit therefore has scaled transfer norm
at most \(e^{L_{\rm real}}<4.3\).
Let \(\Phi_j(s,t)\) denote its physical transfer.
Conjugation by \(S_j\) does not increase either the first-row norm or the
maximum norm of the third column: in the first row the scaling factors
are \(1,\ell_j,\ell_j^2,1\), and in the third column they are
\(\ell_j^2,\ell_j,1,\ell_j^2\).
Thus, with \(e_3=(0,0,1,0)^t\),
\begin{equation}\label{cert:partial-transfer}
 \|(\Phi_j(s,t))_{1,\cdot}\|_{\rm row}\le4.3,\qquad
 \|\Phi_j(s,t)e_3\|_\infty\le4.3
 \quad(s_j\le t\le s\le s_{j+1}).
\end{equation}
Only these two parts of a physical transfer are needed.  Its full norm
need not satisfy the same estimate.

All exponential comparisons just used, and the one in
\eqref{cert:disk-h2}, can be checked by rational arithmetic.
For \(0\le x<82\), let
\[
 E_{80}(x)=\sum_{j=0}^{80}\frac{x^j}{j!}
 +\frac{x^{81}}{81!}\frac1{1-x/82}.
\]
The positive exponential tail has successive ratios at most \(x/82\),
so \(e^x\le E_{80}(x)\).
Substituting the displayed rational values in this formula verifies all
the strict numerical comparisons.  The supplementary
\texttt{verification/arithmetic\_allowances.py} performs these rational
checks, including \eqref{cert:initial-error},
\eqref{cert:product-perturbation}, and the remaining comparisons below.

\subsection{Uniform nonlinear error and closure of the bootstrap}
\label{cert:nonlinear}

Fix \(z_*\in[-R,R]^2\).  On each block compare its actual shot with the
reference solution on that block, and write \(\delta p=p-p_{\rm ref}\).
For the only nonlinear term in the physical third-component equation,
\[
 -\frac{p}{1+p}
 +\frac{p_{\rm ref}}{1+p_{\rm ref}}
 +\frac{\delta p}{(1+p_{\rm ref})^2}
 =\frac{(\delta p)^2}{(1+p_{\rm ref})^2(1+p)}.
\]
It follows that the difference equation, beyond the reference
linearization, has forcing \(e_3\mathcal N\), where
\begin{equation}\label{cert:nonlinear-remainder}
 |\mathcal N|=
 \frac{|\delta p|^2}{|1+p_{\rm ref}|^2|1+p|}.
\end{equation}
The \(d\)-dependence is linear and is already included in the fourth
variation component.

Temporarily assume \(|\delta p|\le14R\) throughout the portions of the
blocks reached so far.  Since \(m_j>.23\) and
\(m_j/(m_j-14R)<1.001\), \eqref{cert:nonlinear-remainder} implies
\[
 |\mathcal N|\le\frac{1.001(14R)^2}{m_j^3}.
\]
The error accumulated by the end of block \(j\), before multiplication
by subsequent complete-block transfers, has zero fourth component.
Its norm is bounded by
\[
 10^{-19}
  +4.3\,\ell_j\frac{1.001(14R)^2}{m_j^3}.
\]
The first term is the reference jump
\eqref{cert:local-errors}; the second follows by variation of constants
and the third-column estimate in \eqref{cert:partial-transfer}.

Define the endpoint error \(e_i\) by
\[
 \begin{pmatrix}p(s_i)\\p'(s_i)\\p''(s_i)\\d\end{pmatrix}
 =\begin{pmatrix}U_i\\d_0\end{pmatrix}
          +T_{i,0}Zz_*+e_i.
\]
Its fourth component is zero.  By \eqref{cert:initial-error},
\eqref{cert:exact-sums}, and the preceding local error estimate,
\begin{align}
 \|e_i\|_\infty
 &\le580000\cdot10^{-19}+8501\cdot(2\cdot10^{-19})
       +4.3\cdot250000\cdot1.001(14R)^2\notag\\
 &=.1055747504R<.11R.
 \label{cert:endpoint-error}
\end{align}
For an interior point of a block, the first-row estimate in
\eqref{cert:partial-transfer}, \eqref{cert:linear-size}, and the forcing
estimate give
\begin{align}
 |\delta p|
 &\le4.3(2.72+.11)R
      +4.3\frac{1.001(14R)^2}{.23^3}\notag\\
 &<12.169035R<14R.
 \label{cert:bootstrap-improvement}
\end{align}
Here replacing the remaining block length by \(1\) is a harmless
overestimate.  This improves the bootstrap strictly both at block starts
and in block interiors, including the reference jumps.
At \(s=1\) it starts from \(\|Z\|<.02\) and
\eqref{cert:initial-error}.  Continuation over the finitely many blocks
therefore proves \eqref{cert:endpoint-error} at every block endpoint.

For clarity, existence is part of this continuation argument.
As long as the bootstrap holds, \(1+p\ge .23-14R>0\).
On the finite interval \(1\le s\le56\), the equations for \(p'\) and
\(p''\) form a linear system with bounded coefficients and bounded forcing
once \(p\) stays in this range.  Gronwall's inequality bounds these
derivatives on every such interval.  Thus no other finite-time escape
can interrupt continuation.  Positivity before \(s=1\) was established
in Subsection~\ref{cert:origin-error}.  This proves the existence and
positivity statements of Proposition~\ref{cert:shoot}.

\subsection{Projection at the matching point}\label{cert:projection}

Let \(\Gamma\) denote the exact real coordinate matrix of
\eqref{prof:coordinates}.
Lemma~\ref{prof:root-bounds} gives
\[
 |\Gamma_{ab}-\Lambda_{ab}|<.001,
 \qquad
 \|\Gamma_{1,\cdot}\|_{\rm row}<.92,\qquad
 \|\Gamma_{a,\cdot}\|_{\rm row}<.68\quad(a=2,3).
\]
Thus \(\|\Gamma-\Lambda\|<.003\).
The polynomial \(P_{\rm a}\) also satisfies
\begin{equation}\label{cert:pa-variation}
 \|P_{\rm a}(d)-P_{\rm a}(d_0)\|_\infty
 \le .01|d-d_0|\le .01\cdot.011R;
\end{equation}
this follows by differentiating the four explicit coefficients in
\eqref{prof:asymptotic-polynomial}, as detailed in the proof of
Lemma~\ref{prof:tail-residual}.

Set \(A_*=U_N-P_{\rm a}(d_0)\).
The endpoint error and transfer error give
\[
 \begin{pmatrix}p\\p'\\p''\end{pmatrix}(56)-P_{\rm a}(d)
 = A_*+Yz_*+\epsilon_*
       -\bigl(P_{\rm a}(d)-P_{\rm a}(d_0)\bigr),
 \qquad
 \|\epsilon_*\|_\infty\le
       \bigl(.11+.001\cdot.02\bigr)R.
\]
In fact \(\epsilon_*\) is the sum of the top three entries of \(e_N\)
and \(((T_{N,0}-\widetilde T_{N,0})Zz_*)_{\mathrm{top}\,3}\).
Apply one row of \(\Gamma\).
The contribution of \(A_*\), using \eqref{cert:finite-matching}, is at
most \((.006+.003\cdot.04)R\).
For the first row the linear contribution is bounded by
\((.116+.003\cdot2.7)R\).
For row \(a=2,3\), after subtracting the corresponding coordinate of
\(z_*\), it is bounded by \((.001+.003\cdot2.7)R\).
The remaining contribution is bounded by the row norm of \(\Gamma\)
times
\[
 \bigl(.11+.001\cdot.02+.01\cdot.011\bigr)R.
\]
Consequently, in units of \(R\), the two required estimates are
\begin{align*}
 |q|/R
 &\le .006+.003\cdot.04+.116+.003\cdot2.7
       +.92(.11+.001\cdot.02+.01\cdot.011)\\
 &=.2315396<.25,\\
 \big\|(\Re\xi,\Im\xi)-z_*\big\|_\infty/R
 &\le .006+.003\cdot.04+.001+.003\cdot2.7
       +.68(.11+.001\cdot.02+.01\cdot.011)\\
 &=.0901084<.13.
\end{align*}
These are exact rational comparisons and establish
\eqref{cert:conclusion}, completing the proof of
Proposition~\ref{cert:shoot}.

\subsection{Executable source and reproducibility}
\label{cert:provenance}

The complete source printed in Appendix~\ref{app:code} is
\texttt{verification/certificate.py}, comprising \(6594\) UTF-8 bytes.
Its SHA-256 digest is
\begin{center}\small\ttfamily
263389fd7217d84d0e5b66869797fc9f7998\\
309d524a95f7ad7d3d0f538ea742
\end{center}
where the line break is for typesetting only.
It uses no imported libraries, floating-point operations, external data,
or platform-dependent numerical routines.  From the paper directory,
the verification command is
\begin{quote}\ttfamily
python3 verification/certificate.py
\end{quote}
with Python assertions enabled (no \texttt{-O} or \texttt{-OO}, and no
\texttt{PYTHONOPTIMIZE} setting that disables assertions).
The program is silent when every assertion succeeds.

The supplementary rational check
\begin{quote}\ttfamily
python3 verification/arithmetic\_allowances.py
\end{quote}
verifies the analytic allowances used above by exact fraction arithmetic;
its decimal report is descriptive only.
The optional script \texttt{verification/tail-check.py} checks the explicit
root and tail estimates in Section~\ref{sec:profile} by exact rational
arithmetic.
The finite integer assertions, the analytic enclosure argument in
Subsections~\ref{cert:origin-error}--\ref{cert:disks}, and the nonlinear
argument in Subsections~\ref{cert:transfer-errors}--\ref{cert:projection}
together prove the finite shooting proposition.
